\section{可解释性、调试与自动化}

\subsection{Q 值可视化与诊断}

在训练过程中，可以通过 heatmap 或 contour 观察 Q 值随动作、状态变化的分布。若某个动作一直得到极高 Q 分数却从未选出，可能表明行为策略实现（如 $\epsilon$-greedy）被 bug 影响。

\begin{knowledgebox}{诊断工具要素}
重要信息包括：Q 值梯度、TD error 曲线、Replay buffer 中 action frequency。可用 TensorBoard、MLFlow 来展示时间序列图。
\end{knowledgebox}

\subsection{训练流水线的自动化}

推荐在训练周期中嵌入自动化检查：每隔若干 epoch 进行 offline evaluation、同步 metrics 到 dashboard、自动计算 sample efficiency。训练 pipeline 还应支持 `checkpoint+resume` 和 `metric gating`，后者在关键指标下降时自动暂停训练。

\begin{importantbox}{Replay buffer 哨兵}
用一个轻量级模型实时计算 buffer 的 coverage 和 reward statistics，一旦发现分布尖锐偏移就发出警报。这种哨兵可以避免在早期就采集到 one-hot 数据导致后续训练崩溃。
\end{importantbox}

\begin{warningbox}{自动化过度依赖历史指标}
如果 gating 只看历史 reward 平均值，则可能忽略 recent decay。建议引入 sliding window version 的指标，将异常提高到 operator 手动确认。
\end{warningbox}

\subsection{本章小结}

可解释性与自动化并行是规模化训练的基础，整理好诊断视图和自动 gating 机制能显著减少重启次数。

